% =====================================================================================
% Chapter 13 · Endogenous exposure — instrumental variables
%
% Built from notebooks/11_endogenous_exposure_iv.ipynb (the body) and
% notebooks/11b_endogenous_exposure_real_data.ipynb (§13.10, the Criteo companion).
% NOT ONE NUMBER IN THIS FILE IS TYPED BY HAND: every figure in the prose is a macro
% from build/macros.tex, generated by cmp.report from the executed notebooks. The only
% literals below are (i) the constants of the data-generating model in §13.2, which are
% *definitions* rather than results; (ii) the decision-rule threshold (0.9) and the F > 10
% convention, which are *conventions* stated in the text, not measurements.
% =====================================================================================
\chapter{Endogenous exposure: instrumental variables}
\label{ch:iv}

\begin{quote}\small
\emph{An ad platform reports that exposed users convert far more --- and it is telling the truth.
On Criteo's incrementality experiment the exposed visit the advertiser \nbElevenBNaive{} percentage
points more often than the unexposed, with a standard error of \nbElevenBNaiveSe{} pp: the tightest
interval anywhere in this chapter, wrapped around a number that misses the
\nbElevenBNFullM-million-row anchor by \nbElevenBNaiveErr{} pp --- an anchor its own interval
excludes. The platform chose whom to expose, and it chose people who were already going to buy. A
randomized lockout that no user could see instruments the exposure, prices it at
\nbElevenBAnchorWald{} pp, and turns a dashboard bid cap of \eur{\nbElevenBFullNaiveCap} per exposed
user into a defensible one of \eur{\nbElevenBCap}. The chapter's second finding is about the
machinery rather than the marketing: on those rows two-stage least squares and the posterior agree
to \nbElevenBAgreePp{} pp --- \nbElevenBAgreeSes{} classical standard errors --- and the posterior's
interval is \nbElevenBWidthRatio$\times$ the classical one. This is the chapter where Bayes buys the
least, and \cref{sec:iv:compare} says so.}
\end{quote}

% =====================================================================================
\section{The decision}
\label{sec:iv:decision}

A demand-side platform sends a quarterly report. Users who saw the campaign's ads convert at a
far higher rate than users who did not; the platform proposes to bid harder, and asks the firm to
raise its cap on the price it will pay per exposed user. The report's number is not a fabrication.
It is a correctly computed difference of means between two groups of real customers.

It is also useless, and for a reason that has nothing to do with arithmetic. Nobody flipped a coin
to decide who saw the ad. The platform did --- and it is paid to expose the people most likely to
convert. It bids more for users who have been browsing the category, who searched the brand last
week, who linger. Those users would have converted at an elevated rate whether or not an ad ever
reached them. The exposed and the unexposed are not two arms of an experiment; they are the output
of a targeting rule that was optimizing for exactly the outcome now being measured.

The consequence is a specific, priceable error. Let $X_i \in \{0,1\}$ record exposure, $Y_i$ the
outcome, and $U_i$ the customer's unobserved purchase intent --- the thing the auction was bidding
on and the analyst never sees. The naive regression of $Y$ on $X$ is consistent for the causal
effect only if exposure is uncorrelated with everything else that drives the outcome, and here it
is correlated with the largest thing that does:
%
\begin{equation}
  \operatorname{Cov}\big(X_i,\ \underbrace{12\,U_i + \varepsilon_i}_{\text{the full error}}\big)
  \;\neq\; 0 .
  \label{eq:iv:endog}
\end{equation}
%
\Cref{eq:iv:endog} is worth reading slowly, because the condition is stated against the
\emph{whole} error term and not against the idiosyncratic noise $\varepsilon$ alone. The noise is
independent of exposure by construction and always will be; it is $U$ --- omitted, unrecorded,
unrecordable --- that leaks. This is what \emph{endogenous} means: the treatment is correlated with
the hidden drivers of the outcome. No amount of controlling for what \emph{was} measured fixes it.
The adjustment toolkit of the earlier chapters --- back-door sets, propensity scores, doubly robust
estimation --- assumes the confounder is in the data. Here, by construction, it is not.

The cost of believing the platform is not abstract. In this chapter's simulated market the naive
number prices an exposure at \eur{\nbElevenNaive} against a truth of \eur{\nbElevenTrue}; buying a
million exposures a quarter at an OLS-calibrated cap authorizes \eur{\nbElevenOverpayM}\,M per
quarter for value that does not exist. On Criteo's real rows the same error runs at
\eur{\nbElevenBPremium} per exposed user, or roughly \eur{\nbElevenBPremiumQuarter} per quarter at
a million exposures. What is needed is a source of variation in exposure that the platform did not
choose --- and a customer's intent could not have caused.

% =====================================================================================
\section{The data-generating model}
\label{sec:iv:dgm}

A causal method cannot be graded on real data, because on real data the answer is missing --- that
is the entire problem. So the estimator is first marked against a world whose answer is known by
construction, and only then trusted on one whose answer is not. \Cref{sec:iv:real} performs the
second half of that bargain on \nbElevenBNFull{} rows of a real advertising experiment; this
section builds the first.

The simulator draws $n = \nbElevenN$ customers. Each carries an \textbf{unobserved} engagement
$U_i$ --- the intent the auction can see and the analyst cannot --- and receives a
\textbf{randomized} encouragement $Z_i$: a serving-priority lottery that makes the ad more likely
to reach them but is not itself an ad. Writing $\sigma(\cdot)$ for the logistic function,
%
\begin{align}
  U_i &\sim \mathcal N(0, 1), \qquad Z_i \sim \text{Bernoulli}(\tfrac12),
        \label{eq:iv:latent}\\
  X_i &\sim \text{Bernoulli}\Big(\sigma\big(0.3 + 1.1\,Z_i + 0.8\,U_i\big)\Big),
        \label{eq:iv:first}\\
  Y_i &= 50 \;+\; 15\,X_i \;+\; 12\,U_i \;+\; \varepsilon_i,
        \qquad \varepsilon_i \sim \mathcal N\!\left(0,\ 6^{2}\right).
        \label{eq:iv:dgp}
\end{align}
%
The planted causal effect of an exposure is \eur{\nbElevenTrue}, and it is the grade every
estimator in this chapter is marked against.

Everything the chapter is about is visible in where $U$ appears. It enters \emph{both}
\cref{eq:iv:first} and \cref{eq:iv:dgp}: $0.8\,U$ pushes a customer toward exposure, $12\,U$ pushes
the same customer toward buying. That double appearance \emph{is} the self-selection, and it is
what makes \cref{eq:iv:endog} bite. Because the model is known, the damage can be computed before
a single regression is run. The omitted-variable-bias formula, specialised to \cref{eq:iv:dgp},
says the OLS slope converges not to 15 but to
%
\begin{equation}
  \hat a_1 \;\xrightarrow{\;p\;}\; 15 \;+\;
  \frac{\operatorname{Cov}\big(X,\ 12U + \varepsilon\big)}{\Var(X)}
  \;=\; 15 \;+\; 12\,\frac{\operatorname{Cov}(X, U)}{\Var(X)},
  \label{eq:iv:ovb}
\end{equation}
%
the second equality holding because $\varepsilon$ is independent structural noise. Even the sign is
knowable in advance: the first stage loads on $U$ with a positive coefficient, so
$\operatorname{Cov}(X,U) > 0$ and OLS \emph{must} come out too high. Evaluated on the simulator's
own (normally invisible) $U$, \cref{eq:iv:ovb} predicts $15 + \nbElevenOvb =
\eur{\nbElevenOlsPredicted}$. The regression, run afterwards, returns \eur{\nbElevenNaive}. Theory
called the number before the data was allowed to speak, and that is the most useful thing a
simulator ever does: on real data the same formula still gives the \emph{sign} of the bias even
when its size cannot be computed.

The instrument $Z$ is the escape. It shifts exposure --- the coefficient $1.1$ in
\cref{eq:iv:first} is not zero --- and it appears nowhere in \cref{eq:iv:dgp}. The slice of exposure
that $Z$ drives is therefore as good as randomly assigned, uncontaminated by $U$. IV isolates that
slice and reads the effect off it. \Cref{fig:iv:dag} states the structure in one picture, and the
important thing about the picture is what is missing from it.

\input{build/floats/nb11_dag}

% =====================================================================================
\section{Identification}
\label{sec:iv:ident}

\subsection*{The estimand, and why it is not the ATE}

Write $X_i(z)$ for the exposure customer $i$ would end up with if the encouragement were
\emph{set} to $z$, and $Y_i(x)$ for the outcome under exposure $x$. Two potential exposures
generate four customer types, and the whole of IV's interpretation lives in that table.

\begin{center}\small
\begin{tabular}{@{}lll@{}}
\toprule
 & $X_i(1) = 1$ & $X_i(1) = 0$ \\
\midrule
$X_i(0) = 1$ & \textbf{Always-taker} --- the loyalist who & \textbf{Defier} --- so annoyed by the \\
             & seeks the ad out; the lottery changes & nudge they block the ad they would \\
             & nothing & otherwise have seen \\[3pt]
$X_i(0) = 0$ & \textbf{Complier} --- the lottery tips them & \textbf{Never-taker} --- ad-blocked,
                                                             barely \\
             & into seeing the ad & online; unreachable either way \\
\bottomrule
\end{tabular}
\end{center}

The object IV identifies is the average effect \emph{among compliers} --- the local average
treatment effect, or LATE \citep{imbens1994, angrist1996}:
%
\begin{equation}
  \beta_{\text{LATE}} \;=\; \E\big[\,Y_i(1) - Y_i(0) \;\big|\; X_i(1) > X_i(0)\,\big].
  \label{eq:iv:late}
\end{equation}
%
This is \textbf{not} the average treatment effect. Always-takers and never-takers are not merely
down-weighted in \cref{eq:iv:late}; they are \emph{absent} from it. And --- the part most often
lost in translation --- the complier population is \textbf{defined by the instrument}.
A different encouragement moves a different set of customers, and so estimates a different
quantity, even on identical data with an identical outcome. Two valid instruments that
disagree are not evidence that one is broken \citep{imbens2014, deaton2010}. Report the
scope, always.

For a budget decision this is the right scope rather than a compromise: a bid buys
\emph{incremental} exposure, so the relevant question is what an exposure does for the customers a
bid can actually move.

\subsection*{The assumptions}

Identification of \cref{eq:iv:late} rests on four conditions. Two of them are arrows the instrument
must \emph{not} have, which is exactly why no dataset can test them.

\begin{assumption}[Relevance]\label{as:iv:rel}
The instrument moves the treatment: $\pi \neq 0$ in the first stage $X = b_0 + \pi Z + u$. This is
the one assumption the data can check, and \cref{sec:iv:classical} checks it with the first-stage
F. It is measured here at $F = \nbElevenFStat$, and \cref{sec:iv:wrong} degrades the instrument on
purpose to show what happens when it fails.
\end{assumption}

\begin{assumption}[Exclusion]\label{as:iv:excl}
The instrument affects the outcome \emph{only} through the treatment. In potential-outcome
notation, writing $Y_i(z, x)$ for the outcome when encouragement is set to $z$ and exposure to $x$,
%
\begin{equation}
  Y_i(z, x) \;=\; Y_i(x) \qquad \text{for all } z \in \{0, 1\}.
  \label{eq:iv:exclusion}
\end{equation}
%
Holding exposure fixed, flipping the lottery changes nothing. \textbf{This is untestable from data
and it is the whole ballgame.} It is defended on design --- a serving-priority lottery sells nothing
by itself; an encouragement \emph{email that shows the product} would violate it flatly --- and it
is \emph{priced} rather than assumed away in \cref{sec:iv:euro}, where a hypothetical direct effect
$\delta$ is subtracted from the reduced form and the $\delta$ that would overturn the decision is
computed.
\end{assumption}

\begin{assumption}[Exogeneity]\label{as:iv:exog}
$Z \perp (U, \varepsilon)$: the instrument is independent of the hidden drivers of both exposure
and outcome. Here it holds \emph{by randomization} --- the lottery is a fair coin. On real data this
is the condition that separates an instrument from a hopeful covariate, and \cref{sec:iv:real}
audits it directly, by checking that twelve user features are balanced across the instrument while
being wildly imbalanced across exposure.
\end{assumption}

\begin{assumption}[Monotonicity / no defiers]\label{as:iv:mono}
$X_i(1) \ge X_i(0)$ for every customer: the encouragement never pushes anyone \emph{away} from
exposure. Untestable, usually plausible for a gentle nudge --- and, as the derivation below shows,
it is precisely what turns a ratio of two experimental contrasts into an interpretable average
effect. In \cref{sec:iv:real} it holds \emph{by construction}: locked-out users cannot be exposed at
all, so a defier is a logical impossibility rather than an assumption.
\end{assumption}

\subsection*{The LATE theorem, in two lines}

Decompose the intent-to-treat effect --- the average outcome difference between encouraged and
unencouraged customers, which is nothing but the reduced form --- over the four types.
Randomization (\cref{as:iv:exog}) makes the two groups comparable. Exclusion
(\cref{as:iv:excl}) says the outcome can only react to $Z$ if exposure changed, which silences
always-takers and never-takers, for whom $X_i(1) = X_i(0)$. Monotonicity (\cref{as:iv:mono})
deletes the defier term. What survives is the compliers:
%
\begin{equation}
  \underbrace{\E[Y \mid Z{=}1] - \E[Y \mid Z{=}0]}_{\text{ITT / reduced form}}
  \;=\; \E\big[\big(Y_i(1) - Y_i(0)\big)\big(X_i(1) - X_i(0)\big)\big]
  \;=\; \Prob(\text{complier}) \cdot \beta_{\text{LATE}} .
  \label{eq:iv:itt}
\end{equation}
%
The identical decomposition applied to exposure itself gives
$\E[X \mid Z{=}1] - \E[X \mid Z{=}0] = \Prob(\text{complier})$: \textbf{the first stage
\emph{is} the complier share}. Dividing \cref{eq:iv:itt} by it isolates the LATE, and the ratio has
a name older than the computer \citep{wright1928}:
%
\begin{equation}
  \hat\beta_{\text{Wald}}
  \;=\; \frac{\E[Y \mid Z{=}1] - \E[Y \mid Z{=}0]}{\E[X \mid Z{=}1] - \E[X \mid Z{=}0]}
  \;=\; \frac{\hat\delta}{\hat\pi}
  \;=\; \frac{\text{what the lottery did to sales}}{\text{what the lottery did to exposure}} .
  \label{eq:iv:wald}
\end{equation}
%
In this market the lottery lifts the exposure rate from \pc{\nbElevenFsLo} to \pc{\nbElevenFsHi}: a
shift of \nbElevenFsShift{} points, which under \cref{as:iv:mono} is the estimated complier share.
About a fifth of these customers are compliers, and they are the only ones IV speaks for.

% =====================================================================================
\section{The classical baseline}
\label{sec:iv:classical}

Before any sampler, the discipline of this book is to ask what a competent analyst produces
\emph{without} one. In this chapter the answer is unusually blunt: three least-squares slopes and a
division. The causal work was done by \crefrange{as:iv:rel}{as:iv:mono}, not by any machinery, and
IV's machinery predates the computer.

\subsection*{Three regressions, naked, in this order}

\begin{equation}
  \underbrace{Y = a_0 + a_1 X}_{\textbf{(1) naive OLS --- biased}}, \qquad
  \underbrace{X = b_0 + \pi Z}_{\textbf{(2) first stage}}, \qquad
  \underbrace{Y = c_0 + \delta Z}_{\textbf{(3) reduced form}} .
  \label{eq:iv:three}
\end{equation}
%
Regression (1) is the number the business would otherwise have used, and \cref{eq:iv:ovb} already
convicted it. Regressions (2) and (3) each regress on the \emph{randomized} lottery, so each is a
clean experiment --- but neither is the answer, because one is denominated in exposures and the
other in euros. Divide, the units cancel into euros \emph{per exposure}, and
\cref{eq:iv:wald} is recovered. \Cref{tab:iv:rungs} runs all four rows.

\input{build/tables/nb11_rungs}

\begin{proposition}[2SLS \emph{is} the Wald ratio]\label{prop:iv:wald}
With a single binary instrument, a single binary endogenous regressor and no controls, two-stage
least squares and $\hat\delta/\hat\pi$ are the same estimator. The notebook asserts this in code
rather than asking for faith: it computes both and checks
$|\hat\delta/\hat\pi - \hat\beta_{\text{2SLS}}| < 10^{-8}$.
\end{proposition}

Which raises the obvious question --- why run 2SLS at all, if the answer is a division? For the two
things a bare ratio cannot carry: a \textbf{standard error} and the \textbf{first-stage F}. Both
turn out to repay the trouble.

\subsection*{The standard error the folklore gets backwards}

Two-stage least squares is often taught as a recipe: regress $X$ on $Z$, keep the fitted $\hat X$,
regress $Y$ on $\hat X$. Followed literally, the recipe gives the right \emph{coefficient} and the
wrong \emph{standard error}, and the received wisdom about the direction of that error is
\textbf{false}. The usual warning is that the hand-rolled second stage comes out too narrow. Here it
comes out too \emph{wide}, by a factor of \nbElevenSeHandRatio (\cref{tab:iv:se}), and the algebra
says it must.

Write $b$ for the 2SLS coefficient. The correct residual is formed against the real regressor,
$e = y - Xb$. The hand-rolled second stage forms it against the fitted one, and the two differ by
exactly the first-stage prediction error:
%
\begin{equation}
  \underbrace{y - \hat X b}_{\text{what the naive fit residualises}}
  \;=\; \underbrace{(y - Xb)}_{e,\ \text{the correct residual}}
      \;+\; \underbrace{(X - \hat X)\,b}_{\text{first-stage prediction error}} .
  \label{eq:iv:resid}
\end{equation}
%
The two terms on the right are orthogonal in sample --- $\hat X$ is the projection of $X$ onto the
instrument space, so $X - \hat X$ is orthogonal to everything the second stage regresses on --- and
therefore the sums of squares simply add:
%
\begin{equation}
  \text{SSR}_{\text{naive}} \;=\; \text{SSR}_{\text{correct}}
    \;+\; b^{2}\!\sum_i (X_i - \hat X_i)^{2}
  \;\;\ge\;\; \text{SSR}_{\text{correct}} .
  \label{eq:iv:ssr}
\end{equation}
%
The extra term is a sum of squares: non-negative, always. The hand-rolled fit therefore charges
itself for first-stage \emph{prediction} error, which is not sampling error in $\hat\beta$ at all,
and reports an interval that is \textbf{conservative}. Here the inflation is
\nbElevenSeHandRatio$\times$: \eur{\nbElevenSeHand} against the correct \eur{\nbElevenSeCorrect}.
That is not a licence to be sloppy --- a wrong standard error is wrong in either direction, and it
will not stay conservative when controls or multiple instruments enter --- but it is worth stating
plainly, because a chapter that repeated the folklore would have been repeating something the
notebook can measure and refute.

\input{build/tables/nb11_se}

The covariance choice needs one line of justification, as everywhere in this book. This is a
cross-section of \nbElevenN{} customers: no panel, no time series, so no clustering and no HAC. But
exposure is binary, so the outcome variance differs between the exposed and unexposed arms by
construction --- the errors are heteroskedastic \emph{by design}. All three regressions of
\cref{eq:iv:three} therefore take HC1 standard errors \citep{white1980, mackinnon1985}. For the IV
estimate itself the homoskedastic 2SLS formula and an HC1 sandwich agree to
\nbElevenSeRobustRatio$\times$ (\cref{tab:iv:se}), so the interval printed above is safe as it
stands.

\subsection*{Relevance, and what the first-stage F is \emph{for}}

For a single instrument the first-stage F is the ordinary F-statistic of regression (2),
%
\begin{equation}
  F \;=\; \frac{R^{2}/1}{(1 - R^{2})/(n - 2)},
  \label{eq:iv:F}
\end{equation}
%
with $R^2$ the share of exposure variance the instrument explains: the instrument's signal in the
first stage relative to its noise. The folklore bar is $F > 10$ \citep{staiger1997, stock2005}, and
the bar is not arbitrary. With one instrument, a standard approximation prices the 2SLS estimator's
residual bias \emph{relative to the OLS bias it was hired to remove} at roughly the reciprocal of
the F:
%
\begin{equation}
  \frac{\E[\hat\beta_{\text{2SLS}}] - \beta}{\E[\hat\beta_{\text{OLS}}] - \beta}
  \;\approx\; \frac{1}{F} .
  \label{eq:iv:oneoverF}
\end{equation}
%
So $F = 10$ means: no more than about a tenth of the OLS bias leaks back in. That is a rule with
teeth rather than a superstition, and it can be cashed out in euros on any particular instrument.
Here OLS is off by \eur{\nbElevenNaiveErr}; at $F = \nbElevenFStat$, \cref{eq:iv:oneoverF} prices
the residual leakage at \eur{\nbElevenFLeak} per exposure. Negligible --- which is precisely why the
small residual bias \cref{sec:iv:valid} finds in the \emph{posterior mean} cannot be blamed on the
instrument. \Cref{sec:iv:wrong} runs the rule backwards and watches the estimator come apart.

\subsection*{The grade, and the guardrail}

Against the planted \eur{\nbElevenTrue}: naive OLS returns \eur{\nbElevenNaive}, off by
\eur{\nbElevenNaiveErr}, and the truth sits \textbf{outside} its 90\,\% interval of
$[\eur{\nbElevenNaiveLo}, \eur{\nbElevenNaiveHi}]$. Two-stage least squares returns
\eur{\nbElevenIvEst}, off by \eur{\nbElevenIvErr}, and the truth sits comfortably \emph{inside}
$[\eur{\nbElevenIvLo}, \eur{\nbElevenIvHi}]$. The gap between them --- \eur{\nbElevenWedge} per
exposure --- \emph{is} the self-selection the instrument removed.

One detail in that paragraph deserves to be pulled out and set in bold, because it is the single
most useful thing in the classical arm. \textbf{The naive interval is the tightest in the chapter
and it misses the truth entirely.} Its standard error is \eur{\nbElevenNaiveSe} --- a half-width of
\eur{\nbElevenNaiveHalf} against 2SLS's \eur{\nbElevenIvHalf} --- and it is wrapped around a number
that is \eur{\nbElevenNaiveErr} wrong. More data would not have saved it; more customers would only
have shrunk that interval further around the same wrong number. \textbf{Precision is not accuracy,
and a confident interval on a biased estimator is a confident lie.}

What this classical apparatus cannot produce, at any level of effort, is $\Prob(\beta > c)$: the
probability that an exposure is worth more than the price the firm would have to bid for it. A
confidence interval is a property of the \emph{procedure} --- 90\,\% of intervals built this way
would cover the truth across repeated samples --- and this particular interval either covers the
truth or it does not. \Cref{sec:iv:euro}'s rule is nothing but that probability. That is the gap the
next section exists to fill --- and, as \cref{sec:iv:compare} will insist, it is the \emph{only} gap
the next section fills.

% =====================================================================================
\section{The Bayesian model}
\label{sec:iv:bayes}

The Bayesian IV changes the identification by not one comma: same instrument, same four
assumptions, same complier LATE. It changes the apparatus for expressing uncertainty about that
number, and it estimates one object the classical route structurally cannot hand over --- the
endogeneity itself.

The model fits the first stage and the outcome equation \emph{jointly}, allowing their errors to be
correlated:
%
\begin{align}
  \begin{pmatrix} X_i \\ Y_i \end{pmatrix} &\;\sim\;
    \mathcal N_2\!\left(
      \begin{pmatrix} b_0 + \pi Z_i \\ \beta_0 + \beta X_i \end{pmatrix},\ \Sigma \right),
    \qquad
    \Sigma = \begin{pmatrix}
       \sigma_X^{2} & \rho\,\sigma_X\sigma_Y \\
       \rho\,\sigma_X\sigma_Y & \sigma_Y^{2} \end{pmatrix},
    \label{eq:iv:joint}\\
  b_0,\ \pi,\ \beta_0,\ \beta &\;\sim\; \mathcal N\!\left(0,\ 50^{2}\right),
    \qquad
  \Sigma \;\sim\; \text{LKJCholeskyCov}(\eta = 2).
    \label{eq:iv:priors}
\end{align}
%
The off-diagonal $\rho$ \textbf{is} the endogeneity: the correlation between the two equations'
errors that the hidden $U$ induces and that OLS ignores. The LKJ prior \citep{lkj2009} is a
standard weakly-informative distribution over correlation matrices; at $\eta = 2$ it pulls $\rho$
gently toward zero, so the data has to argue for endogeneity rather than have it assumed. Inference
is by the No-U-Turn sampler \citep{hoffman2014} in PyMC \citep{salvatier2016}, via CausalPy:
\nbElevenNDraws{} post-warmup draws, $\hat R = \nbElevenRhat$, minimum bulk effective sample size
\nbElevenEss{} \citep{vehtari2021}, \nbElevenDivergences{} divergent transitions. The chains have
converged.

The fit returns a posterior mean of \eur{\nbElevenBayesMean}, with a 90\,\% credible interval of
$[\eur{\nbElevenBayesLo}, \eur{\nbElevenBayesHi}]$ --- against 2SLS's \eur{\nbElevenIvEst} and
$[\eur{\nbElevenIvLo}, \eur{\nbElevenIvHi}]$ (\cref{fig:iv:bias}). The two agree to
\eur{\nbElevenEstDiff}, and their interval widths --- \eur{\nbElevenBayesWidth} against
\eur{\nbElevenIvWidth} --- are the same to within a rounding error. That is by design, and it is
the first thing \cref{sec:iv:compare} will hold the sampler to.

\input{build/floats/nb11_bias}

\subsection*{The prior, priced rather than warned about}

A book that told the reader ``choose weakly-informative priors'' and moved on would be asserting
exactly the kind of thing this book exists to measure. So the notebook \emph{fits} the alternative.
CausalPy's library default centres the coefficient priors on the 2SLS point estimate with a
standard deviation of one. Same data, same sampler, only the priors differ:

\begin{itemize}[leftmargin=*, itemsep=.25em]
  \item \textbf{Library default:} posterior mean \eur{\nbElevenPriorDefMean}, 90\,\% interval
        $[\eur{\nbElevenPriorDefLo}, \eur{\nbElevenPriorDefHi}]$ --- width
        \eur{\nbElevenPriorDefWidth}. The planted \eur{\nbElevenTrue} lies \textbf{outside} it.
  \item \textbf{$\mathcal N(0, 50^2)$, used here:} width \eur{\nbElevenPriorWeakWidth}, and the
        truth is inside.
\end{itemize}

The default band is \nbElevenPriorNarrowFactor$\times$ too narrow, and the independent yardstick is
not a matter of taste: the classical Wald interval, which knows nothing about any prior, is
\eur{\nbElevenIvWidth} wide. A prior centred on the point estimate with $\sigma = 1$ does not let
the data speak. It hard-codes the answer and then reports false confidence in it --- and the failure
is not a Bayesian failure but a defaulting failure, of the kind that ships silently because nobody
fitted the alternative. This chapter fitted it.

\subsection*{The one genuinely new object: $\rho$}

Two-stage least squares \emph{removes} the endogeneity by projection. It never names it, and it has
no $\rho$ to report at any width. The joint model estimates it:
$\rho = \nbElevenRho$, with a 90\,\% interval of $[\nbElevenRhoLo, \nbElevenRhoHi]$. That posterior
is the Bayesian image of exactly the $\operatorname{Cov}(X, U)$ that \cref{eq:iv:ovb} priced at
\eur{\nbElevenOvb} --- a positive $\rho$ \emph{is} the reason the naive number sat above the causal
one. On real data, where there is no answer key, a credibly positive $\rho$ is evidence that the
dashboard number \emph{was} inflated: evidence the classical arm cannot produce.

% =====================================================================================
\section{Diagnostics and validation}
\label{sec:iv:valid}

\subsection*{The check the model is designed to fail}

Exposure is 0/1. \Cref{eq:iv:joint} models it with a Gaussian. Something has to give, and the
posterior predictive check is where it gives (\cref{fig:iv:ppc}). Replicate datasets simulated from
the fitted joint model reproduce the \emph{sales} margin well --- replicate standard deviation
\eur{\nbElevenPpcRepSd} against an observed \eur{\nbElevenPpcObsSd}, with a faint under-dispersion
--- which is what licenses reading $\beta$ off the outcome equation. But \pc{\nbElevenPpcOutPct} of
the replicated \emph{exposures} land outside $[0, 1]$: exposures of $-0.3$ and $1.4$, values no
customer can have.

\input{build/floats/nb11_ppc}

This is not a bug to be fixed. It is the linear-probability approximation made visible, and running
the check is how a modelling note stops being a footnote and becomes a measurement. It has two
consequences the chapter honours throughout: the exposure equation's \emph{point estimate} and
$\rho$ may be read, its posterior \emph{width} may not; and the small bias the next subsection
measures has a named suspect.

\subsection*{Recovery across fresh samples}

One sample can land off by luck. The honest question is whether the estimator is centred on the
truth over repeated samples and whether its interval covers at the stated rate. Refitting on
\nbElevenNSeeds{} fresh datasets: the posterior mean averages \eur{\nbElevenSeedMean} against a
planted \eur{\nbElevenTrue} --- a bias of \eur{\nbElevenSeedBias}, or \pc{\nbElevenSeedBiasPct} ---
with a seed-to-seed standard deviation of \eur{\nbElevenSeedSd}, and its 90\,\% interval covers the
truth in \nbElevenSeedCovN{} of \nbElevenNSeeds{} samples (\pc{\nbElevenSeedCovPct} against a
90\,\% target: modest under-coverage, consistent with that bias).

\subsection*{Whose fault is the bias --- the identification's, or the model's?}

This is the diagnostic that pays for the simulator, and it is a clean experiment. The Wald ratio
\emph{is} the IV identification with nothing else attached: two least-squares slopes and a division,
no prior, no likelihood, no sampler. Re-run it on fresh draws at growing $n$ and whatever bias it
does \emph{not} show cannot belong to the identification.

\input{build/tables/nb11_waldn}

It shows none (\cref{tab:iv:waldn}). The closed-form Wald stays pinned on the planted truth at
every sample size --- a bias of \eur{\nbElevenWaldBiasBig} at $n = 10{,}000$ --- while its spread
shrinks like $1/\sqrt{n}$. The verdict is therefore precise, and it does not flatter the sampler:
\textbf{the identification is clean, and the posterior mean's residual bias is the Bayesian model's
own} --- the Gaussian first stage \cref{fig:iv:ppc} just caught in the act, plus the ordinary
skewness of a ratio posterior at modest draw counts. \Cref{eq:iv:oneoverF} already priced the pure
IV leakage at \eur{\nbElevenFLeak}, an order of magnitude smaller. \textbf{The Bayesian machinery
introduced a modelling error that the two \code{lstsq} calls of \cref{sec:iv:classical} could not
have made}, because they never modelled $X$ at all.

% =====================================================================================
\section{Point estimate versus posterior}
\label{sec:iv:compare}

\input{build/tables/nb11_compare}

\subsection*{On the number, they agree --- and Bayes does not win}

\Cref{tab:iv:compare} puts all three estimators on the estimand of \cref{eq:iv:late}, on the same
\nbElevenN{} customers, under the same four assumptions. Only the apparatus differs.

The two causal rows land within \eur{\nbElevenEstDiff} of each other and both sit slightly above
the planted truth; both intervals cover it. On this seed the \textbf{classical estimator is the
closer of the two} --- \eur{\nbElevenIvErr} against the posterior's \eur{\nbElevenBayesErr} --- and
\cref{sec:iv:valid} has already explained why: the residual in the classical arm is ordinary
sampling noise on one draw, while the posterior mean's extra drift is a modelling error of its own
making. Say it plainly: \textbf{Bayes did not win the point estimate.} Nor did it win the width. The
half-widths are \eur{\nbElevenIvHalf} (2SLS) and \eur{\nbElevenBayesHalf} (posterior) --- the same
scale, and deliberately so, since the weakly-informative priors of \cref{eq:iv:priors} were chosen
precisely so that the data and not the prior would set the width. The posterior bought no precision,
and was never supposed to.

This is the ordinary case in this book, and it has a name. With a well-specified likelihood and
enough data the posterior concentrates where the maximum-likelihood estimate does, with the same
sampling variance --- the Bernstein--von Mises phenomenon \citep{vandervaart1998, gelman2013}. It
will appear again in \cref{sec:iv:real}, in a form so exact it can be read straight off a print-out.

\subsection*{What the posterior did buy}

One thing --- and the last column of \cref{tab:iv:compare} is it. $\Prob(\beta > c)$ exists in
exactly one of the three rows. The classical entry reads \emph{not defined}, and that is a literal
statement rather than a rhetorical concession. A confidence interval attaches no probability to a
hypothesis about a fixed parameter.

And --- this is the part practitioners get wrong --- \textbf{bootstrapping the Wald ratio would not
fix it}. The share of bootstrap resamples landing above $c$ is a frequency over hypothetical
\emph{datasets}; it is not a belief about $\beta$. Resampling harder does not convert a statement
about a procedure into a statement about a parameter. The quantity simply does not exist in the
classical vocabulary. It exists in a posterior:
%
\begin{equation}
  \Prob(\beta > c) \;=\;
  \frac{1}{S}\sum_{s=1}^{S} \mathbf 1\big[\beta^{(s)} > c\big],
  \qquad \beta^{(s)} \sim p(\beta \mid \text{data}),
  \label{eq:iv:pgt}
\end{equation}
%
and \cref{sec:iv:euro}'s bid-cap rule consumes nothing else. Second, $\rho$: an interval on the very
bias that 2SLS silently projects away.

\subsection*{The ledger}

Not a better point estimate; the classical one is nearer the truth. Not a tighter interval; the
widths match by design, and the \emph{tightest} interval in the chapter belongs to the biased OLS
row. Not protection from model error --- the Bayesian arm \emph{created} one, in the
linear-probability first stage of \cref{fig:iv:ppc}. Not the weak-instrument diagnostic: the F is a
classical object, and the posterior would have gone on producing a beautifully smooth, entirely
fictitious credible band as F fell through 10. What it bought is the \emph{decision}: a probability
statement about the effect, propagated coherently into euros, plus a read-out of the endogeneity
that caused the problem in the first place.

That is the whole of the win. \Cref{sec:iv:real} will show that on thirteen million real rows it is
even less than that --- and that it is still, just, enough.

% =====================================================================================
\section{The decision, in euros}
\label{sec:iv:euro}

The plan is to buy exposures at \eur{\nbElevenCost} each. The decision rule is a stated convention
applied to the posterior of \cref{eq:iv:pgt}:
%
\begin{equation}
  \text{bid up to } c \quad\text{while}\quad \Prob(\beta > c) \;\ge\; 0.9,
  \qquad\Longleftrightarrow\qquad
  c^{*} \;=\; q_{0.10}(\beta),
  \label{eq:iv:rule}
\end{equation}
%
where $q_p(\beta)$ is the posterior $p$-quantile. The equivalence on the right is worth pausing on,
because it makes the cap's status unambiguous: sweeping a hypothetical cost against a \emph{fixed}
posterior just walks that posterior's quantile function, so the budget cap \textbf{is} a posterior
quantile denominated in euros. It is a \textbf{90\,\% probability bar} --- a statement about
$\beta$ --- and it is emphatically \emph{not} a 90\,\% confidence bar. No confidence interval in
\cref{tab:iv:compare} can be read this way, and no bootstrap of them can be made to.

At \eur{\nbElevenCost} the posterior gives $\Prob(\text{pays}) = \nbElevenPPays$ and a net value of
\eur{\nbElevenNet} per exposure. The probability is saturated rather than assumed: the entire
90\,\% credible interval $[\eur{\nbElevenBayesLo}, \eur{\nbElevenBayesHi}]$ sits above the cost
line. The cap of \cref{eq:iv:rule} lands at \eur{\nbElevenCap} per exposure, the coin-flip price at
\eur{\nbElevenCoinflip}, so the current plan carries \eur{\nbElevenHeadroom} of genuine headroom
(\cref{fig:iv:euro}). And the counterfactual mistake has a price: an OLS-calibrated cap would have
sanctioned paying up to \eur{\nbElevenNaive} for something causally worth \eur{\nbElevenTrue}. At
\nbElevenVolume{} exposures a quarter, that authorizes roughly \eur{\nbElevenOverpayM}\,M per
quarter --- \eur{\nbElevenOverpayMLo}\,M to \eur{\nbElevenOverpayMHi}\,M across the IV interval ---
for value that does not exist.

\input{build/floats/nb11_euro}

\subsection*{Pricing the untestable assumption}

\Cref{as:iv:excl} cannot be tested. It can, however, be \emph{priced}, and that is the difference
between an assumption and a wish. Suppose the lottery had a direct effect $\delta$ on sales --- the
encouragement email carries a discount code, say, or shows the product. Then the honest reduced form
is $\hat\delta_{\text{obs}} - \delta$, and the corrected estimate is
%
\begin{equation}
  \beta(\delta) \;=\; \frac{\hat\delta_{\text{obs}} - \delta}{\hat\pi},
  \qquad\text{so}\qquad
  \beta(\delta) > c \;\iff\; \delta \;<\; \delta^{*} = \hat\delta_{\text{obs}} - c\,\hat\pi .
  \label{eq:iv:excl_stress}
\end{equation}
%
The buy call survives until $\delta \approx \eur{\nbElevenBreakDelta}$ ---
\pc{\nbElevenBreakDeltaPct} of the entire \eur{\nbElevenReduced} reduced form would have to bypass
exposure to overturn it (\cref{fig:iv:euro}, right). For a content-free serving lottery that is a
large leak, and the decision is robust. For an encouragement email that doubles as an
advertisement, it is not remotely a large leak, and the design is broken before the data is
collected. \textbf{Design the instrument so that \cref{eq:iv:excl_stress} is a formality}; do not
hope to rescue a leaky one afterwards.

% =====================================================================================
\section{What can go wrong}
\label{sec:iv:wrong}

\subsection*{Weak instruments, demonstrated rather than deplored}

Everything above rested on a strong first stage. \Cref{eq:iv:oneoverF} says what that buys; the
honest way to show what its absence costs is to take it away. The encouragement coefficient of
\cref{eq:iv:first} is turned down from \nbElevenWeakGammaTop{} to nearly nothing and the
closed-form Wald ratio recomputed on \nbElevenWeakNDraws{} fresh datasets at each setting --- the
failure being demonstrated belongs to the \emph{identification}, so no sampler is involved and the
whole sweep costs seconds (\cref{tab:iv:weak}, \cref{fig:iv:weak}).

\input{build/tables/nb11_weak}
\input{build/floats/nb11_weakfig}

The median holds up longer than the spread does, which is itself the trap: an analyst who looked
only at the point estimate would see nothing alarming until well past the point of no return. At
$F \approx \nbElevenWeakFWorst$ a single draw's 90\,\% range spans \eur{\nbElevenWeakLo} to
\eur{\nbElevenWeakHi} --- far wider than the \eur{\nbElevenNaiveErr} of OLS bias the instrument was
hired to remove --- and the median has drifted back up toward OLS. The mechanism is
\cref{eq:iv:oneoverF} read right to left: with a nearly irrelevant instrument the Wald denominator
$\hat\pi$ is sampling noise around zero, so the ratio inherits OLS's bias \emph{and} an enormous
variance \citep{bound1995, staiger1997}.

That is the precise sense in which \textbf{a weak instrument is worse than none}. OLS is wrong by a
known-ish amount with a tight interval. A weak IV can be wrong by multiples of the effect in either
direction \emph{while wearing the costume of a causal estimate}. Check the F before believing any IV
number; if it is small, walk away rather than report the ratio. Three refinements belong in a
practitioner's toolkit and are named here rather than derived: on real, heteroskedastic data prefer
the Montiel Olea--Pflueger \emph{effective} F, which corrects the threshold for non-constant
first-stage noise \citep{olea2013}; with multiple endogenous regressors the rank test of
\citet{kleibergen2006} replaces the scalar F; and $F > 10$ should be read as necessary, not
sufficient \citep{andrews2019}.

\subsection*{The failure modes that no diagnostic will catch}

\Cref{as:iv:excl} and \cref{as:iv:mono} are untestable, and calling them ``assumptions'' rather than
``diagnostics'' is not modesty --- it is the accurate description. An instrument that leaks a direct
effect produces a biased estimate that no F, no posterior predictive check and no placebo can flag.
\Cref{eq:iv:excl_stress} bounds the damage; institutional knowledge of how the instrument was
generated is the only thing that discharges it.

And \cref{eq:iv:late} is a scope restriction, not a footnote. IV speaks for compliers. Always-takers
--- the loyalists who seek the ad out --- and never-takers are silent in the estimate. A number
earned on the complier margin licenses \emph{a bid on incremental exposure}; it does not license a
claim about what the campaign does to the customer base as a whole, and it does not transfer to a
different instrument that moves a different set of people.

% =====================================================================================
\section{On real data}
\label{sec:iv:real}

Everything above was graded against a truth that was planted. The objection --- that recovering a
number someone put there is not proof of anything --- deserves a real answer, and there is one.

Criteo's advertising incrementality tests publish \nbElevenBNFull{} real users
\citep{diemert2018}. The dataset is a gift for IV, because the test infrastructure randomizes
\textbf{eligibility}: a random 15\,\% of users are locked out of the campaign entirely, and the rest
are eligible to be targeted. Eligibility is a textbook encouragement instrument for self-selected
exposure --- what is randomized is the \emph{nudge}, not the ad. Of the eligible,
\pc{\nbElevenBPExposedEligible} actually end up seeing an ad, chosen by the auction on signals the
analyst cannot observe. Of the locked out, \pc{\nbElevenBPExposedLocked} do.

That last number --- an exact zero, a fact of the serving infrastructure rather than a sample
accident --- earns \cref{as:iv:mono} for free. \textbf{One-sided noncompliance}: a defier would be a
user who gets exposed \emph{only because} they were locked out, which is not implausible but
impossible. And it does more than that. Returning to \cref{eq:iv:itt}, the term
$\E[(Y(1) - Y(0))X \mid Z{=}0]$ is exactly zero when locked-out users are never exposed, so the
ratio collapses to
%
\begin{equation}
  \hat\beta_{\text{Wald}}
  \;=\; \frac{\E\big[(Y(1) - Y(0))\,X \mid Z{=}1\big]}{\Prob(X{=}1 \mid Z{=}1)}
  \;=\; \E\big[\,Y(1) - Y(0) \mid X = 1\,\big]
  \;=\; \text{ATT of the exposed},
  \label{eq:iv:bloom}
\end{equation}
%
the conditioning on $Z{=}1$ dropping because every exposed user is an eligible one
\citep{bloom1984}. The compliers here are \emph{literally} the exposed --- which is exactly the
population a bid cap is about.

\subsection*{The referee: a design-based anchor, not a planted truth}

Nothing was planted, so nothing can be marked against a truth. What replaces it is the
\textbf{full-file anchor}: the ITT, the first stage and their ratio computed \emph{exactly},
design-based, on all \nbElevenBNFull{} rows, where sampling noise is negligible at the precision
that matters. It is not a truth; it is the number this same identification strategy converges to
when one stops subsampling. On the full file the ITT is \nbElevenBAnchorItt{} pp and the first stage
\nbElevenBAnchorFirst, giving a Wald LATE of \textbf{\nbElevenBAnchorWald{} pp} per exposure --- and
a full-file naive as-treated gap of \nbElevenBAnchorNaive{} pp. Every estimate below is graded
against it. Substituting a full-population design-based estimate for ground truth is the honest
real-data move, and it is more than most applied projects ever get.

\subsection*{The confounding, photographed}

Here is what a simulator can never show (\cref{fig:iv:smd}). Twelve anonymized user features ship
with the data. Across the \emph{randomized} eligibility flag their standardized mean differences are
flat --- a maximum of \nbElevenBSmdZMax{} --- which is \cref{as:iv:exog}, audited rather than
asserted. Across \emph{actual exposure} the same features are skewed by as much as
\nbElevenBSmdDMax{} standard deviations. The exposed \emph{were different people before any ad
ran}: the unobserved $U$ of \cref{eq:iv:dgp}, measured.

\input{build/floats/nb11b_smd}

And this is only $U$'s observable shadow. The auction also bid on signals that are not in the file
--- which is precisely the regime in which covariate adjustment surrenders and an instrument is
required. No regression on those twelve features could have rescued the naive number.

\subsection*{The classical rungs, on real rows}

\input{build/tables/nb11b_rungs}

\Cref{tab:iv:realrungs} runs \cref{eq:iv:three} on a \nbElevenBNSub-row working subsample. Four
readings, each a different number a business could act on.

\begin{enumerate}[leftmargin=*, itemsep=.35em]
  \item \textbf{The naive as-treated gap is the dashboard's number, and it is wrong by a third.}
        It comes to \nbElevenBNaive{} pp against the anchor's \nbElevenBAnchorWald{} pp, off by
        \nbElevenBNaiveErr{} pp --- and its 90\,\% interval, at
        $[\nbElevenBNaiveLo, \nbElevenBNaiveHi]$ pp, \emph{excludes} the anchor. It is also the
        tightest interval in the section (SE \nbElevenBNaiveSe{} pp): thirteen million more rows
        would only shrink it further around the same wrong number. \Cref{sec:iv:classical}'s lesson,
        reprinted on real data.
  \item \textbf{The ITT is honest and useless on its own.} At \nbElevenBReduced{} pp it is unbiased
        for the \emph{intent} --- randomization guarantees it, and it lands within a hair of the
        full-file \nbElevenBAnchorItt{} pp. But it is the effect of \emph{eligibility}, averaged over
        the roughly 96\,\% of eligible users the auction never reached. As a price \emph{per
        exposure} it understates by the reciprocal of the compliance rate: a factor of
        $1/\hat\pi \approx \nbElevenBDilutionFactor$. Reporting the ITT as ``the ad's effect'' is the
        mirror-image error to reporting the naive gap. One is diluted, the other inflated. Only their
        ratio is denominated in the unit a bid is priced in.
  \item \textbf{The Wald ratio undoes exactly that dilution --- and it is 2SLS.} It returns
        \nbElevenBWald{} pp, \nbElevenBWaldErr{} pp from the anchor, which sits comfortably inside
        its interval. \Cref{prop:iv:wald} holds here too, and the notebook asserts it in code again.
        The wedge the instrument removed, on the full file: \nbElevenBAnchorWedge{} pp of the
        \nbElevenBAnchorNaive{} pp dashboard number --- \pc{\nbElevenBAnchorWedgeShare} of it --- was
        intent the exposed users already had. Quoted against the causal number instead of the
        dashboard's, the same wedge means the dashboard runs \pc{\nbElevenBAnchorInflation} too big.
        Two denominators, one wedge; name the one you are dividing by.
  \item \textbf{The F is a certificate, not a worry.} $F = \nbElevenBFStat$, which is
        \nbElevenBFOverBar$\times$ the $F > 10$ bar. Eligibility is randomized and mechanically
        enforced by the ad server, not scraped from behaviour; this is about as strong an instrument
        as a real dataset ever hands anyone. \Cref{eq:iv:oneoverF}, cashed out: the naive number is
        off by \nbElevenBNaiveErr{} pp, so the residual OLS-bias leakage into 2SLS is
        $\approx \nbElevenBFLeak$ pp. A rounding error. This is the regime IV was designed for and
        rarely finds, and the diagnostic is here to \emph{certify}, not to worry.
\end{enumerate}

The interval for row (4) is the delta method, because the estimator is a ratio and both its parts
are noisy:
%
\begin{equation}
  \Var\big(\hat\beta\big) \;\approx\;
    \frac{1}{\hat\pi^{2}}\Var\big(\hat\tau\big)
  \;+\; \frac{\hat\beta^{2}}{\hat\pi^{2}}\Var\big(\hat\pi\big)
  \;-\; \frac{2\hat\beta}{\hat\pi^{2}}\operatorname{Cov}\big(\hat\tau,\ \hat\pi\big),
  \label{eq:iv:delta}
\end{equation}
%
the three terms being the numerator's noise, the denominator's noise, and their correlation
(positive: a slice that happens to draw high-intent users gets both more exposure \emph{and} more
visits). Because the enormous F pins $\hat\pi$ down so precisely --- its relative noise is exactly
$1/\sqrt{F}$ --- the last two terms nearly vanish, and the one-term shortcut
$\text{SE} \approx \text{SE}(\hat\tau)/\hat\pi$ comes within \pc{\nbElevenBSeShortGap} of the full
expression: \nbElevenBSeShort{} pp against \nbElevenBSeDelta{} pp. That approximation is
\emph{audited}, not asserted, and it earns its keep in the fold check below. And 2SLS's own algebra
returns \nbElevenBSeTwosls{} pp --- a ratio of \nbElevenBSeTwoslsRatio{} to the hand-propagated
variance. On a saturated binary design they are the same object, computed two ways.

\subsection*{Fresh samples, without a simulator}

\Cref{sec:iv:valid}'s multi-seed recovery cannot be run here: there is only one world. But
\nbElevenBNFullM{} million rows afford fresh \emph{samples}. Cutting the working subsample into
\nbElevenBNFolds{} disjoint folds of \nbElevenBFoldRows{} rows and refitting the same estimator on
each, the fold Walds average \nbElevenBFoldMean{} pp with a standard deviation of
\nbElevenBFoldSd{} pp around the \nbElevenBAnchorWald{} pp anchor, and their 90\,\% intervals cover
it in \pc{\nbElevenBFoldCover} of folds (\cref{fig:iv:folds}). That spread is the honest price of
subsampling a \pc{\nbElevenBPExposedEligible} exposure rate --- and the reason the anchor is
computed on the whole file rather than on a slice of it.

\input{build/floats/nb11b_folds}
\input{build/floats/nb11b_anchor}

\subsection*{Bernstein--von Mises, in a print-out}

\input{build/tables/nb11b_compare}

Now the head-to-head, and it is the reason this chapter exists in the form it does.
\Cref{tab:iv:realcompare} re-runs the classical arm on \emph{exactly} the \nbElevenBNFit{} rows the
sampler saw, so the comparison is apparatus against apparatus and not $n$ against $n$. Both are
scored against the full-file anchor. Set expectations before reading it: at a hundred thousand rows
the likelihood swamps any weakly-informative prior, and Bernstein--von Mises says the posterior
should be \emph{approximately the sampling distribution of the classical estimator} --- same centre,
same width. If Bayes bought precision here, something would be wrong.

It did not, and the agreement is closer than ``approximately'' suggests. The two causal rows land
\nbElevenBAgreePp{} pp apart --- \nbElevenBAgreeSes{} classical standard errors --- and their
interval half-widths are \nbElevenBHalfClassical{} pp (delta method) and \nbElevenBHalfBayes{} pp
(posterior), a ratio of \nbElevenBWidthRatio$\times$. \textbf{The posterior standard deviation and
the delta-method standard error are estimating the same number, and they agree.} This is not a happy
accident; it is a theorem \citep{vandervaart1998} showing up in a print-out. The sampler spent
roughly twenty minutes rediscovering, by simulation, a division that took microseconds.

Two-stage least squares is \emph{nominally} the closer of the two to the anchor:
\nbElevenBFitWaldErr{} pp against the posterior's \nbElevenBPostErr{} pp. That gap is
\nbElevenBErrGapSes{} standard errors. \textbf{It is noise, not skill, and this chapter declines to
call a winner on it.} The one row that is decisively wrong is the same one as ever --- the naive
as-treated gap, \nbElevenBFitNaiveErr{} pp from the anchor, carrying the tightest interval in the
table.

(The contrast with \cref{sec:iv:compare} is instructive. On \nbElevenN{} simulated rows the
classical arm was \emph{genuinely} the closer of the two, because the Gaussian first stage of
\cref{fig:iv:ppc} was committing a real linear-probability modelling error the two \code{lstsq} calls
could not make. That error has not gone away here. It has been \emph{drowned}: at Criteo's $n$, model
error of that size is invisible next to nothing at all.)

\subsection*{The honest verdict --- this is where Bayes buys the least}

So say it plainly, in the words the notebook uses rather than in softer ones. \textbf{At this $n$,
Bayes buys nothing on estimation.} Not accuracy: the errors are indistinguishable. Not precision: the
credible interval \emph{is} the confidence interval, at \nbElevenBWidthRatio$\times$ its
width. Not the instrument diagnostic: the F is a classical object, and the posterior would have gone
on producing a beautifully smooth, entirely fictitious credible band as F fell through 10. Not
robustness: the linear-probability approximation is the \emph{Bayesian} arm's liability, not the
classical one's. Anyone who says the Bayesian machinery \emph{rescued} this analysis is selling
something.

What it does buy is the same one thing, and it is not nothing: the decision. The
$\Prob(v\beta > c)$ column of \cref{tab:iv:realcompare} exists in one arm and reads \emph{not
defined} in the other, and \cref{eq:iv:pgt}'s bootstrap objection applies here in full --- a
frequency over hypothetical datasets is not a belief about $\beta$.

\begin{remark}[The lesson, stated so it survives being quoted]
At Criteo scale the Bayesian layer is a \textbf{convenience, not a rescue}. It does not improve the
number; it changes what one is licensed to \emph{say} about the number, and it propagates that
licence coherently into euros. Where Bayes genuinely earns its compute is the opposite regime ---
thin data, hierarchical structure, real prior information, a model that must be built rather than
divided. Here the only honest reason to fit it is the decision rule of \cref{eq:iv:rule}, and
\textbf{if the firm did not need $\Prob(\text{effect} > \text{cost})$, the classical baseline alone
would have been a complete and defensible analysis.}
\end{remark}

\subsection*{The dollar decision, and the cap the dashboard would have set}

At an assumed \eur{\nbElevenBValuePerVisit} per incremental visit, an exposure to the users this
auction actually reaches is worth about \eur{\nbElevenBAnchorCap} (the anchor's
\nbElevenBAnchorWald{} pp). At the planned \eur{\nbElevenBCost} per exposed user,
$\Prob(\text{pays}) = \nbElevenBPPays$ and the net value is \eur{\nbElevenBNet}. The cap of
\cref{eq:iv:rule} --- again, a \textbf{90\,\% probability bar}, a posterior quantile, and not a
confidence bar --- lands at \eur{\nbElevenBCap} per exposed user, with a coin-flip price of
\eur{\nbElevenBCoinflip} (\cref{fig:iv:realeuro}).

\input{build/floats/nb11b_euro}

The attribution dashboard would have set that cap at \eur{\nbElevenBFullNaiveCap}. The difference,
\eur{\nbElevenBPremium} per exposure, is a pure targeting premium paid for intent the firm already
owned: at \nbElevenBExposuresQuarter{} exposures a quarter, roughly
\eur{\nbElevenBPremiumQuarter} every quarter. And the untestable assumption is priced here too, by
\cref{eq:iv:excl_stress}: the buy call survives an eligibility side-channel of up to
$\delta^{*} = \nbElevenBDeltaStar$ pp --- \pc{\nbElevenBBreakDeltaPct} of the entire
\nbElevenBReduced{} pp reduced form would have to bypass exposure to overturn it. The closed form
and the grid scan agree in front of the reader (\nbElevenBDeltaStar{} pp against
\nbElevenBBreakDelta{} pp), which is the only reason to print both.

Three scope lines before any of this reaches a rate card. It is the ATT of the
\pc{\nbElevenBPExposedEligible} of users the auction actually reaches: it prices \emph{these}
exposures and does not license doubling reach, because the next users the auction would find are, by
its own logic, worse prospects. The \eur{\nbElevenBValuePerVisit} visit value and the
\eur{\nbElevenBConvMargin} margin per conversion are finance inputs, swept but assumed --- the
secondary conversion rung puts an exposure at \eur{\nbElevenBConvValue} on the full-file anchor's
\nbElevenBAnchorConvWald{} pp, and at a \pc{\nbElevenBBaseConv} base rate it is honestly noisy.
And the estimate is a blend across several incrementality tests: fine for a benchmark cap, to be
blurred into one campaign's plan with care.

\begin{remark}[What the experiment bought]
Notice what the randomized lockout bought. It made \cref{as:iv:exog} true by construction,
\cref{as:iv:mono} true by construction, and \cref{as:iv:excl} defensible in a sentence --- a silent
lockout does nothing to a user who never learns of it. Every hard problem in this chapter was solved
\emph{in the design}, before any data was analysed. That is the deeper reason to run incrementality
tests at all: not to produce a number, but to make the number identifiable \citep{lewis2015,
gordon2019}.
\end{remark}

% =====================================================================================
\section{Takeaways}
\label{sec:iv:take}

\begin{enumerate}[leftmargin=*, itemsep=.45em]
  \item \textbf{Endogeneity is a correlation with the \emph{full} error, and adjustment cannot
        touch it.} \Cref{eq:iv:endog} is stated against $12U + \varepsilon$, not against
        $\varepsilon$: the leak is the unobserved intent, and controlling for what was measured is
        powerless against a confounder that was not. An instrument does not adjust for $U$. It finds
        a slice of exposure that $U$ could not have caused.

  \item \textbf{IV identifies a LATE, and the instrument defines the population.} \Cref{eq:iv:late}
        is the effect on \emph{compliers} --- \nbElevenFsShift{} points of these customers, the ones
        the encouragement moves. Always-takers and never-takers are absent from the estimate, not
        merely down-weighted. A different instrument therefore estimates a different quantity, and
        two valid instruments disagreeing is not a contradiction. Report the scope.

  \item \textbf{The first-stage F is a precondition, and \cref{eq:iv:oneoverF} tells you what it
        buys.} At $F = \nbElevenFStat$ the residual OLS-bias leakage is \eur{\nbElevenFLeak} per
        exposure; on Criteo, at $F = \nbElevenBFStat$ --- \nbElevenBFOverBar$\times$ the bar --- it
        is \nbElevenBFLeak{} pp. In both cases the diagnostic \emph{certifies}. Below 10 it does the
        opposite: \cref{tab:iv:weak} shows a single draw's 90\,\% range spanning \eur{\nbElevenWeakLo}
        to \eur{\nbElevenWeakHi} at $F \approx \nbElevenWeakFWorst$, with the median drifting back
        toward OLS. A weak instrument is worse than none.

  \item \textbf{The hand-rolled 2SLS standard error is too WIDE, not too narrow.} \Cref{eq:iv:ssr}
        is the proof and \cref{tab:iv:se} the measurement: \nbElevenSeHandRatio$\times$ the correct
        one. Regressing $Y$ on $\hat X$ charges the estimate for first-stage prediction error that is
        not sampling error at all. The folklore says the opposite; the folklore is wrong.

  \item \textbf{Precision is not accuracy.} The tightest interval in this chapter belongs, in both
        the simulation and the real data, to the biased estimator: a standard error of
        \eur{\nbElevenNaiveSe} on a number \eur{\nbElevenNaiveErr} wrong, and \nbElevenBNaiveSe{} pp
        on one \nbElevenBNaiveErr{} pp wrong --- around an anchor its own interval excludes. More
        data shrinks that interval around the same wrong number, forever.

  \item \textbf{Classical and Bayesian agree, and on Criteo they \emph{coincide}.} On
        \nbElevenBNFit{} rows the posterior and 2SLS land \nbElevenBAgreePp{} pp apart ---
        \nbElevenBAgreeSes{} standard errors --- with interval widths in a ratio of
        \nbElevenBWidthRatio$\times$. That is Bernstein--von Mises in a print-out: the posterior
        \emph{is} the sampling distribution the delta method wrote in closed form. Which of the two
        sits nearer the anchor differs by \nbElevenBErrGapSes{} standard errors, and this chapter
        refuses to call a winner on that.

  \item \textbf{This is the chapter where Bayes buys the least --- and the book says so.} It bought
        no accuracy, no precision, no diagnostic, no robustness; it \emph{introduced} a
        linear-probability modelling error the classical arm could not make. What it bought is
        $\Prob(\beta > c)$, which the classical column prints as \emph{not defined} --- and
        bootstrapping the Wald would not fix that, because a frequency over datasets is not a belief
        about $\beta$. The euro cap of \cref{eq:iv:rule} is a \textbf{posterior quantile}, a 90\,\%
        \emph{probability} bar and never a confidence bar. \textbf{At Criteo scale the Bayesian layer
        is a convenience, not a rescue: absent the need for $\Prob(\text{effect} > \text{cost})$, the
        classical baseline alone would have been a complete and defensible analysis.}

  \item \textbf{The exclusion restriction is untestable, and that is why you price it.}
        \Cref{eq:iv:excl_stress} bounds the leak that would flip the call --- \eur{\nbElevenBreakDelta}
        of a \eur{\nbElevenReduced} reduced form in the simulation, \nbElevenBDeltaStar{} pp of a
        \nbElevenBReduced{} pp one on Criteo. Defend it on design, not on data: a content-free
        serving lottery is an instrument, and an encouragement email that shows the product is not.
\end{enumerate}

\notebooknote{%
  \Crefrange{sec:iv:dgm}{sec:iv:wrong} are \code{notebooks/11\_endogenous\_exposure\_iv.ipynb}
  (simulator \code{cmp.dgp.iv\_ad\_exposure}, seed \nbElevenSeed, $n = \nbElevenN$);
  \cref{sec:iv:real} is \code{notebooks/11b\_endogenous\_exposure\_real\_data.ipynb}, which loads
  Criteo's published incrementality panel \citep{diemert2018} directly and computes its
  \nbElevenBNFull-row anchor with \code{cmp.data.criteo\_full\_anchor}. Both run in a fast teaching
  mode (\code{CMP\_FAST=1}, short chains) and a full mode (\code{CMP\_FAST=0}); every number in this
  chapter comes from the full run, emitted by \code{cmp.report} and injected here as a macro. The
  Bayesian IV is CausalPy's joint-Gaussian instrumental-variable model, fitted with PyMC; the
  classical arm is \code{cmp.classical.iv\_2sls}. \citet{angrist2009} remains the standard applied
  reference for everything in this chapter, and \citet{imbens2014} the standard survey.}
